%% ====================================================================
%%  Neural Computation (MIT Press) submission manuscript  --  v2
%%  GENERATED from main_v2.tex by make_nc_v2.py. Do not edit by hand:
%%  edit main_v2.tex, re-run the converter, then re-run
%%  verify/verify_paper_numbers.py.
%%  - 12pt, 1 3/8 in margins, author-date (APA) citations, expanded tables
%%  - Draft mode (double-spaced, full width) as required for submission
%%  Build: pdflatex main_nc_v2.tex  (twice)
%% ====================================================================
\documentclass[12pt]{article}

\usepackage[margin=1.375in,top=1.375in,bottom=1.375in]{geometry}
\usepackage{amsmath,amssymb,mathtools}
\usepackage{booktabs}
\usepackage{graphicx}
\usepackage[utf8]{inputenc}
\usepackage[T1]{fontenc}
\usepackage{xcolor}
\usepackage[colorlinks=true,urlcolor=blue,citecolor=blue,linkcolor=blue]{hyperref}
\usepackage{parskip}
\usepackage{setspace}
\usepackage{enumitem}
\usepackage{siunitx}
\usepackage{float}

% Neural Computation submission: draft = double-spaced, full width
\doublespacing

% Macros used by the (inlined) tables, mirroring main_v2.tex
\newcommand{\sub}[1]{\ensuremath{\mathit{#1}}}
\newcommand{\RMSE}{\ensuremath{\mathrm{RMSE}}}
\newcommand{\MI}{\ensuremath{\mathrm{MI}}}
\newcommand{\TE}{\ensuremath{\mathrm{TE}}}
\newcommand{\dfloor}{\ensuremath{\Delta\RMSE_{\mathrm{floor}}}}
\newcommand{\palt}{0.75\,u}

\newenvironment{hangparas}{\begin{list}{}{\setlength{\leftmargin}{1.5em}\setlength{\itemindent}{-1.5em}\setlength{\itemsep}{0pt}\setlength{\parsep}{0pt}}}{\end{list}}

\title{A falsification protocol for simulated embodied spiking substrates: Separating substrate-carried information from readout-channel information}
\author{Alan Daleth Hern\'andez Barreto}
\date{}
\begin{document}
\maketitle

\begin{abstract}
\noindent
When a population of spiking neurons is read out to control a simulated
body, ``the substrate computes'' is often claimed from a decoded error
or a mutual-information value. We argue that such a claim is not
falsifiable unless the substrate is swapped for a \emph{matched dead
substrate} that preserves the single-neuron marginal firing statistics
(mean, variance, silent fraction) while being statistically independent
of the task, and unless the readout is itself controlled. We build a
minimal simulated embodied hub ($2{,}000$ neurons, fixed sparse wiring,
one calibrated decision map) and run a battery of gates: (B) substrate
swap under a fixed readout; (B2) a \emph{paired population null suite}
(per-neuron circular shift, block shuffle, identity swap, task shuffle)
with a false-positive control and a fixed-weight readout-invariance
test; (C) readout-channel ablation; (D) cross-validated ridge readout of
the \emph{same} spikes; (D\,bis) a small nonlinear (MLP) readout on the
same spikes; (E) memory-demanding tasks; and (F) a recurrence ablation.
The matched dead substrate is task-independent by construction and stays
at the null floor under every readout (one-sided $p=1/64$ per profile,
$6/6$ seeds; under the paired null suite, $216/216$ pooled cells at
$p\le0.05$ for each of the four null types), with zero false positives
in the dead-drive control ($0/24$) and a fixed-weight
readout-invariance result that survives at the resolution of the null
sample ($24/24$, $p=0.0476$). The living substrates do not uniformly
survive. LIF carries the task through a cross-validated linear readout
($\Delta\RMSE_{\mathrm{floor}}=+0.094$ vs.\ the mean-predictor floor;
$\MI$ significant in $24/30$ runs) and a small nonlinear readout does
not change that materially ($+0.105$); a seed-level paired test at
$n=24$ seeds is decisive for LIF ($24/24$ seeds same direction,
$p\approx10^{-5}$). The apparent Izhikevich tracking seen under the fixed
calibrated readout is largely a \emph{readout-channel artifact}:
randomizing the readout projection collapses both living substrates to
the dead floor ($3/3$ seeds), a cross-validated linear readout of the
same IZH spikes leaves RMSE essentially at the floor
($\Delta\RMSE_{\mathrm{floor}}=-0.008$), and a nonlinear readout recovers
only a small amount ($+0.012$, $24/24$ seeds same direction). None of the
substrates shows linearly decodable history in the present-window
population state (Gate E negative), so memory claims are not certified
here. The ranking is not architecture-bound: a second hub
($N=4000$, $p_{\mathrm{re}}=0.02$) reproduces it (LIF $+0.104$ vs.\ IZH
$+0.004$ over the dead null). We conclude: (i) matched dead-substrate
controls for the
\emph{material} of a substrate are cheap, decisive, and should precede
scoreboard comparisons; (ii) readout \emph{robustness} is a stronger
falsification test than a single scoreboard, and we reserve the word
``invariance'' for what our experiments cannot show; and (iii) open-loop
diagnosis cannot certify memory, which must be demonstrated in the closed
loop.
\end{abstract}


\section{Introduction}\label{sec:intro}

Physical and simulated substrates for neural computation---in-vitro
cultures (Kagan et al., 2022), neuromorphic silicon
(Maass, 1997; Lukosevicius & Jaeger, 2009; Wringe et al., 2025), or toy spiking
hubs---are routinely evaluated by coupling them to a task and reporting
a score: a decoded error, a correlation, a bit rate. Two distinct claims
are usually conflated under that score. The first is
\emph{substrate-carried information}: spike patterns produced by the
substrate contain task-relevant statistics. The second is
\emph{readout-channel information}: the fixed mapping used to translate
spikes into a decision happens to recover (or inject) task statistics. A
scoreboard alone cannot tell them apart, and the recent discussion
around ``intelligence in a dish''
(Kagan et al., 2022; Balci et al., 2023; Kagan et al., 2023; Smirnova et al., 2023) shows that this conflation has consequences far
beyond toy models.

The contribution of this note is a \emph{falsification protocol} that
separates the two. It rests on three ingredients, all cheap:

\begin{itemize}[nosep]
\item \textbf{A matched dead substrate} (Gate B). The wiring, bias,
drive and readout are kept fixed; only the per-neuron dynamics change.
The dead substrate fires Poisson processes with per-neuron rates
matched to the LIF marginal histogram, but is statistically independent
of the task. Any substrate that cannot beat this null under a fixed
readout is not falsifiably carrying the task.
\item \textbf{A paired population null suite} (Gate B2). A null built
from a single i.i.d.\ Poisson process matches only the single-neuron
marginal and discards temporal structure; a stronger, still cheap
control is to build \emph{sample-matched} surrogates of the substrate's
own spike trains (per-neuron circular shift, block shuffle, identity
swap, task shuffle) and test the observed statistic against them
(Theiler et al., 1992). The task-shuffle null preserves the real count matrix
verbatim---every per-neuron marginal, every cross-neuron correlation and
the full population temporal structure---and destroys only the pairing
of task signal to windows; it is the strictest population-matched
surrogate we can build. We also add a dead-drive
false-positive control and a fixed-weight readout-invariance test.
\item \textbf{Readout controls} (Gates C, D, D\,bis and F). Gate C
randomizes the readout projection into the fixed decision map; Gate D
replaces the hand-tuned readout with a cross-validated ridge readout
trained on the \emph{same} spikes; Gate D\,bis adds a small nonlinear
(MLP) residual readout to bound how much linearity hides; Gate F
ablates recurrence. Task information that remains recoverable after
\emph{all} controls is evidence the substrate activity itself, rather
than the calibrated interface alone, carries the signal. Throughout we
reserve stronger language (``invariance'') for what the experiments
cannot show, and use \emph{readout robustness} for what they do.
\end{itemize}

Gate E then asks the harder question---whether present-window population
activity carries task \emph{history}---which is the minimum a memory
claim needs. Throughout, the replication unit is the seed. Gates B/D/E
pool $6$ seeds $\times$ $5$ profiles; Gate C uses $3$ seeds; Gate D is
additionally repeated at \emph{$n=24$ seeds} with an exact seed-level
paired permutation test; Gates B2 uses the $24$ seeds $\times$ $10$
profiles with $K=50$ paired nulls per cell.

\section{Related work}\label{sec:related}

Six literatures touch the same question from different angles.

\textbf{(1) Embodied neural computation.} ``Intelligence without
representation'' (Brooks, 1991) and morphology-computation continua
(Pfeifer & Bongard, 2006) established that behavior is jointly determined
by body, wiring and control; more recent work formalizes that for
neuromorphic-body systems (D'Angelo et al., 2026; Bartolozzi et al., 2022).
These works score systems, but do not isolate the substrate's material
from its interface. Our hub is deliberately a \emph{simulated} body; the
protocol is what transfers, not the embodiment.

\textbf{(2) In-vitro substrates and the DishBrain debate.} Kagan et
al.\ reported task-relevant responses of cultured cortical neurons in a
closed loop (Kagan et al., 2022); the follow-up correspondence
(Balci et al., 2023; Kagan et al., 2023) turned precisely on whether
the observed signal is attributable to the biology or to the
measurement/readout interface, and the broader organoid-intelligence
program (Smirnova et al., 2023; Zhou et al., 2026) inherits the same
methodological question. Our protocol is a minimal, fully deterministic
enactment of that critique.

\textbf{(3) Spiking and reservoir computing.} Spiking neuron models
(Maass, 1997; Izhikevich, 2003; Izhikevich, 2004; Gerstner & Kistler, 2002) and liquid
state machines / echo-state networks (Maass et al., 2002; Jaeger, 2001; Lukosevicius & Jaeger, 2009) show that random recurrent populations
can support surprisingly capable online regression. The standard
evaluation there is already readout-shaped (train a readout, report
performance); recent critiques of reservoir-computing benchmarks
(Wringe et al., 2025; Yan et al., 2024) make explicit that task scores alone are
hard to compare and easy to over-read. What is usually missing is the
\emph{dead-material control}---the same readout trained on matched
independent Poisson spikes (Lukosevicius & Jaeger, 2009)---and paired
sample-matched surrogates of the substrate's own activity.

\textbf{(4) Neuromorphic benchmarks.} NeuroBench
(Yik et al., 2025) and related efforts define task suites, metrics
and complexity accounting for neuromorphic algorithms and systems, and
embodied extensions are emerging (D'Angelo et al., 2026). Benchmark
suites standardize \emph{what} to score; they do not, by themselves,
supply the substrate-\emph{material} null that our protocol requires.

\textbf{(5) Population decoding and information measures.} Neural
coding theory provides the tools: population decoding with regularized
linear readouts (Georgopoulos et al., 1986; Quiroga & Panzeri, 2009; Paninski et al., 2004; Hoerl & Kennard, 1970), mutual information and its bias correction
(Shannon, 1948; Kraskov et al., 2004; Panzeri et al., 2007), transfer entropy with
surrogate-based significance (Schreiber, 2000; Walker & Newhall, 2018), and
nonlinear readouts (kernel methods, MLPs, population models)
(Paninski et al., 2004; Steinmetz et al., 2019). Our Gate D uses exactly this
machinery; the dead-substrate null is the novel anchor against which
those measures should be read.

\textbf{(6) Null models and statistical rigor.} Surrogate-data methods
(Theiler et al., 1992), permutation statistics (Ernst, 2004),
pseudo-replication warnings (Hurlbert, 1984), and closed-loop
validation practice in brain-machine interfaces
(Carmena et al., 2003; Nicolelis, 2003) all argue that a structurally matched
null is the difference between an effect and an artifact. Gates~B2 and
C are surrogate tests applied to the \emph{substrate and readout
channel} rather than to a single scalar series.

\section{Methods}\label{sec:methods}

\subsection{Embodied hub and wiring}\label{sec:hub}

$N=2{,}000$ neurons, $N_{\mathrm{win}}=200$ control windows of
$C_W=20$ integration steps ($\SI{10}{ms}$ bins, $\SI{0.5}{ms}$
integration). The input projection
$\mathbf{g}\in\mathbb{R}^{N\times1}$ is standard-normal;
recurrence $\mathbf{W}$ is sparse with connection probability
$p_{\mathrm{re}}=0.01$ and entries drawn
$\mathcal{N}(0,1/\sqrt{N p_{\mathrm{re}}})$ (variance-normalized so the
network operates near stability); the readout projection is
$\mathbf{w}_r=\mathbf{g}$ (channel aligned with the input projection),
with a per-substrate scale $k_{\mathrm{ro}}$ (below). The decision map
is the calibrated one of the reference system
(Zhou et al., 2026),
$\hat{x}=\mathrm{clip}((\theta_{\mathrm{sig}}-0.0615)/0.277,\ 0,\ 0.75)$,
clamped to $\SI{0.75}{m/s}$. Five control profiles
($\mathtt{baseline},\mathtt{multi\_step},\mathtt{sine},\mathtt{descend},
\mathtt{pulse}$), six seeds $\{1,7,13,29,55,91\}$ for Gates B/C/D/E,
and Gate B2 additionally uses five profiles of a broader battery
($\mathtt{ramp},\mathtt{perturb},\mathtt{step\_pulse},\mathtt{chirp},
\mathtt{noise}$). Gate~D is repeated with $n=24$ seeds
(Section~\ref{sec:seeds24}).

\subsection{Substrates and the matched dead null}\label{sec:subs}

Three dynamics on identical wiring, bias, drive and readout constants:

\begin{itemize}[nosep]
\item \textbf{LIF}: leaky integrate-and-fire, $\tau=\SI{20}{ms}$,
threshold $v_{\mathrm{th}}=1$, reset $0$, refractory $\SI{2}{ms}$.
\item \textbf{Izhikevich:} regular spiking parameters [``RS'', (0.02,
0.2, $-65$, $8$)] integrated with the $\SI{0.5}{ms}$ step; adaptation
dimension $u$ and $a=0.02$, $b=0.2$, $c=-65$, $d=8$.
\item \textbf{Matched dead null (Poisson):} per-neuron events drawn
independently each integration step from Bernoulli trials with
probability $r_i$; the $r_i$ are estimated from LIF's per-neuron
firing in a constant-input ($u=0.5$) calibration run
(Theiler et al., 1992), so the marginal per-neuron rates match. Table
\ref{tab:hist} reports the match quality explicitly; the null is
\emph{statistically independent} of $u$ by construction.
\end{itemize}

The key design point is that the match is at the level of the
\emph{single-neuron marginal rate histogram}---mean $\mu_r$, variance,
and the fraction of near-silent neurons---not at the level of any
temporal or population structure. This is what makes Gate~B a
\emph{substrate} test rather than a firing-rate test. Table
\ref{tab:hist} (from \texttt{gate\_b\_summary.json}) gives
$\mu_r=\SI{8.09}{Hz}$, SD $\SI{15.30}{Hz}$, Fano
$\sigma^2/\mu=28.9$ (overdispersed vs.\ a homogeneous Poisson) and
$74\%$ neurons with ~zero estimated rate.

\begin{table}[H]
\centering
\caption{Match quality of the dead null to the LIF marginal:
per-neuron mean rate, SD, Fano factor $\sigma^2/\mu$, and the silent
fraction (neuron with $\hat r_i\approx0$). The null draws Bernoulli
events at the matched per-neuron probabilities, independent of $u$.}
\label{tab:hist}
% TABLA marginal emparejada: nulos Pois LIF-matched (mean/SD/Fano).
% desde gate_b_summary.json calibration.poisson_rate
\begin{tabular}{lrrrr}
\toprule
statistic & LIF marginal ($r_i$) & matched IID Poisson & $N=2000$ neurons & $\mathrm{Fano}$ \\
\midrule
per-neuron mean rate (Hz) & 8.091 & 8.091 & both & -- \\
per-neuron StdDev (Hz) & 15.296 & $\sqrt{\mu}$ only & -- & 28.916 \\
fraction silent ($r_i=0$) & 0.740 & $e^{-\mu}$ expected & 1480/2000 & -- \\
\bottomrule
\end{tabular}
% nota: Fano>1 => marginal sobredisperso vs Poisson homogeneo;
% la coincidencia es exacta en tasas marginales (no en estructura
% temporal), que es lo que hace que el nulo sea 'matched'.
\end{table}

\subsection{Calibration without leakage}\label{sec:calib}

$k_{\mathrm{ro}}$ is fixed per substrate from a \emph{constant-input}
calibration run ($u=0.5$, seed $7$): $k_{\mathrm{ro}}=0.2/
\langle \mathbf{w}_r\!\cdot \mathrm{rate}\rangle$ sets the decision
series at the middle of the map's active range. Crucially, calibration
never sees a \emph{task-varying} input, so no task information leaks
into the readout constant. The same $k_{\mathrm{ro}}$ survives Gate~C:
the randomized projection is rescaled to the same norm
($\|\mathbf{w}_r^{\mathrm{rand}}\|=\|\mathbf{g}\|$), so the constant
remains comparable. The Poisson null uses the LIF-derived $k_{\mathrm{ro}}$
plus a $k_{\mathrm{ro}}^{\mathrm{poisson}}$ from its own constant-input
run, again $0.2/\langle\cdot\rangle$.

\subsection{Gate B2: paired population null suite}\label{sec:gateb2}

Gate~B's dead null is i.i.d.\ and therefore destroys any temporal
structure. To test the substrate against its own activity, for each
profile and seed we build $K=50$ sample-matched surrogates of the
recorded counts and recompute the same readout statistic:

\begin{itemize}[nosep]
\item \texttt{circshift}---per-neuron circular shift of the window
sequence (preserves each neuron's marginal and autocorrelation on the
torus);
\item \texttt{blockshuffle}---shuffle contiguous blocks of windows
(block length $B=5$), destroying long-range task alignment while
keeping short-range structure;
\item \texttt{identityswap}---exchange the window order across neurons
(preserves the exact per-window population envelope, destroying
neuron--task alignment).
\end{itemize}

The observed statistic is compared to the $K$ surrogate values with a
one-sided rank test; a cell is significant at $\alpha=0.05$. As a
false-positive control we rerun the whole pipeline with
\texttt{drive\_gain}=0 (no task signal) on three profiles
($\mathtt{sine}$, \texttt{noise}, \texttt{ramp}) $\times$ $6$ seeds
$\times$ $2$ substrates: no cell may be significant. Finally, a
\emph{fixed-weight readout-invariance} test refits nothing: the ridge
weights learned on the real spikes are frozen and applied unchanged to
real and surrogate counts, so the test measures whether recoverable
information survives the surrogate and not whether a fresh model can
adapt to it.

\subsection{Gate C: readout-channel control}\label{sec:gatec}

For each seed $s$, draw an independent zero-mean normal vector
$\mathbf{w}^{\mathrm{rand}}_s$, center it, and rescale to
$\|\mathbf{g}\|$ (RNG seed $4242$). Keep wiring, bias, drive and the
fixed decision map; only the readout projection changes. If apparent
tracking under the aligned channel survives randomization, it is
substrate-independent; if it collapses to the dead floor, the apparent
tracking depended on the readout alignment (whether the substrate
activity \emph{also} contains readout-recoverable task information is
asked separately by Gate D).

\subsection{Decoders (Gates D and D\,bis)}\label{sec:decoders}

On the \emph{same} recorded spike trains:
\begin{enumerate}[nosep]
\item \textbf{canon}--the fixed calibrated channel
$k_{\mathrm{ro}}\,\mathbf{w}_r\!\cdot(\mathrm{counts}/C_W)$ +
\texttt{decide} (Gate~B reference);
\item \textbf{ridge}--cross-validated L2-regularized linear readout of
population counts with $13$ lags on a log-grid in
$[10^{-5},10]$, fit on interleaved train/val windows (every phase appears
in every fold) solved via the Woodbury identity and verified against a
direct solve (max $|\Delta|<4\times10^{-10}$; \texttt{verify\_ridge.py});
\item \textbf{meanpd}--constant-predictor floor
($\mathrm{var}$-around-run-mean);
\item \textbf{passthru}--$\palt$ directly ($\RMSE=0$ by construction,
the no-substrate ceiling);
\item \textbf{mlp} (Gate D\,bis)--the fitted ridge readout plus a small
one-hidden-layer tanh network ($2000\!\to\!32\!\to\!1$, L2, gradient
clipping) trained on the ridge's train residual, with its own penalty
selected on the validation split. This \emph{linear-skip residual} form
starts exactly at the optimal linear readout, so a nonlinear gain is
possible only if it generalizes on held-out windows; a larger net
overfits (the split has only $\approx66$ train windows), which is itself
a data-budget observation for a hub of this size.
\end{enumerate}
We report the \emph{RMSE improvement over the mean-predictor floor},
\[
\dfloor = \RMSE_{\mathrm{meanpd}}-\RMSE_{\mathrm{decoder}},
\]
taking the constant predictor as the zero-information baseline:
positive values mean the decoder beats a constant. We deliberately avoid
an information-unit name for $\dfloor$ because $-\RMSE$ is a loss, not
an entropy; the reader should not read it as a bit or nat quantity.
When comparing a living substrate to the dead null we use the
difference-in-differences $\dfloor(\text{living})-\dfloor(\text{null})$
so that any overfitting of the same decoder on independent spikes
cancels.

\subsection{Statistics and pseudo-replication}\label{sec:stats}

The unit of replication is the \emph{seed} (wiring + initial state).
Gate B/E pool $6$ seeds; Gate D pools $6$ seeds $\times$ $5$ profiles
$=30$ runs, and is repeated at $n=24$ seeds. We report both per-seed
outcomes (e.g.\ ``significant MI in $24/30$ runs'') and pooled means,
and we flag seed-level counts rather than inflating $n$. This follows
standard pseudo-replication guidance (Hurlbert, 1984): the $200$
windows within a run are not independent evidence. The seed-level test
is an exact paired sign-permutation over the $24$ seeds
(Section~\ref{sec:seeds24}); with $n=24$ its one-sided resolution floor
is $\approx10^{-5}$ at $10^5$ permutations.

\subsection{Eager statistics for Gate B}\label{sec:gatebstat}

For the smallest $n$, we report the exact paired permutation over the
$6$ seeds: a one-sided $p=1/64=0.0156$ and two-sided $p=2/64=0.03125$,
with $6/6$ seeds in the same direction. We state the resolution
explicitly rather than quoting deeper $p$-values than the design
supports.

\noindent
\textbf{Experimental design at a glance.}
\begin{table}[H]
\centering
\caption{The gates: what is changed, what is held fixed, and the
question each gate answers. Gate~F (recurrence ablation) was a remark in
v1 and is promoted here to a gate.}
\label{tab:design}
\small
\begin{tabular}{@{}p{5.0em}p{11em}p{11em}p{16.5em}@{}}
\toprule
\textbf{Gate} & \textbf{Manipulation} & \textbf{Held fixed} &
\textbf{Question} \\
\midrule
B (substrate
swap) & dynamics $\to$ Poisson
null with matched marginals
& wiring, bias, drive,
readout $\mathbf{w}_r=\mathbf{g}$, decision map
& Does a task-independent
dead substrate reproduce
the score? \\
B2 (paired
null suite) & counts $\to$ sample-matched
surrogate (circshift / blockshuffle /
identityswap), $K=50$
& wiring, drive, decision map;
plus a dead-drive FP control and a
fixed-weight readout-invariance test
& Does the score survive
surrogates of the substrate's
\emph{own} activity, and is the
test well-calibrated? \\
C (channel
ablation) & readout projection
$\to$ per-seed random $\mathbf{w}_r^{\mathrm{rand}}$
& wiring, drive, spikes,
decision map
& Does the apparent tracking
survive breaking the aligned
channel? \\
D (readout
re-learn) & decoder $\to$ cross-validated ridge on
the \emph{same} spikes
& spikes, task
& Do the spike trains themselves
contain readout-recoverable task
information? \\
D\,bis (nonlinear
readout) & ridge $+$ small tanh MLP residual
on the same spikes
& spikes, task
& How much does the linear
readout hide? \\
E (memory) & task $\to$ history-dependent
(\texttt{lag}, \texttt{integ},
\texttt{stair}); decoder on
present-window state
& spikes, dynamics, readout
& Is task \emph{history} linearly
decodable from the present
population state? \\
F (recurrence) & recurrence $\mathbf{W}\to0$
(nohub), same readout
& drive, readout, decision map
& Does the tested recurrence
contribute at all? \\
G (architecture) & hub $\to$ architecture~B
($N=4000$, $p_{\mathrm{re}}=0.02$),
same substrates \& protocol
& substrates, profiles, seeds,
decoders, decision map
& Is the substrate ranking bound
to the one hub size/density? \\
\bottomrule
\end{tabular}
\end{table}
Gates B/B2/C/D/D\,bis/F decompose one score into (substrate material)
vs.\ (interface); Gate E sets the bar any memory claim must clear; Gate~G
checks the ranking is not architecture-bound. All
are run on the \emph{same} seeds and profiles (Gate~G on a second hub).

\section{Results}\label{sec:results}

\subsection{Gate B: substrate swap under a fixed readout}\label{sec:gateb}
\begin{table}[H]
\centering
\caption{Gate B: per-profile $\widehat{\RMSE}$, $\MI$ (nats), $\rho$
and $\TE$-significance under the fixed calibrated readout (5 profiles,
6 seeds). ``msd'' is the across-seed mean.}
\label{tab:gateb}
% TABLA Gateway B: sustrato intercambiable, decode fijo (5 perfiles x 6 seeds)
% Generada desde gate_b_summary.json. No editar a mano.
\begin{tabular}{llrrrrr}
\toprule
profile & substrate & $\widehat{\mathrm{RMSE}}$ & $\mathrm{MI}$ & $\rho$ & $\mathrm{TE}$ & $p_{\mathrm{TE}}$ \\
\midrule
\texttt{baseline} & LIF & 0.238 & 0.000 & 0.000 & 5e-04 & 0.253 \\
\texttt{baseline} & Izhikevich & 0.254 & 0.000 & 0.000 & 4e-05 & 0.400 \\
\texttt{baseline} & Poisson (dead null) & 0.399 & 0.000 & 0.000 & 0e+00 & 0.833 \\
\texttt{descend} & LIF & 0.167 & 1.034 & 0.933 & 0e+00 & 0.027 \\
\texttt{descend} & Izhikevich & 0.196 & 0.627 & 0.892 & 0e+00 & 0.000 \\
\texttt{descend} & Poisson (dead null) & 0.382 & 0.002 & -0.006 & 0e+00 & 0.833 \\
\texttt{multi\_step} & LIF & 0.131 & 0.681 & 0.994 & 0e+00 & 0.000 \\
\texttt{multi\_step} & Izhikevich & 0.215 & 0.444 & 0.822 & 0e+00 & 0.000 \\
\texttt{multi\_step} & Poisson (dead null) & 0.432 & 0.000 & 0.012 & 0e+00 & 0.833 \\
\texttt{pulse} & LIF & 0.124 & 0.574 & 0.958 & 0e+00 & 0.000 \\
\texttt{pulse} & Izhikevich & 0.294 & 0.241 & 0.570 & 0e+00 & 0.000 \\
\texttt{pulse} & Poisson (dead null) & 0.372 & 0.000 & 0.004 & 0e+00 & 0.833 \\
\texttt{sine} & LIF & 0.179 & 0.966 & 0.938 & 0e+00 & 0.027 \\
\texttt{sine} & Izhikevich & 0.236 & 0.614 & 0.879 & 0e+00 & 0.000 \\
\texttt{sine} & Poisson (dead null) & 0.375 & 0.000 & -0.012 & 0e+00 & 0.833 \\
\bottomrule
\end{tabular}
\end{table}
Under the fixed calibrated readout, LIF and IZH both beat the matched
dead null on every task-varying profile (Table~\ref{tab:gateb}).
The dead null stays at $\MI\approx0$ and $\rho\approx0$; its $\TE$ is
zero under the trial-shuffle estimator at every profile and seed. On
$\mathtt{baseline}$ (constant input) \emph{no} substrate carries
information, as expected. Taken alone, Gate~B would read as ``the
neural substrates compute''---this is precisely the reading that Gates
B2, C and D correct.

\subsection{Gate B2: paired null suite, calibration and invariance}\label{sec:gateb2res}
\begin{table}[H]
\centering
\caption{Gate B2: per substrate and surrogate type, the number of
task-varying cells (profiles $\times$ seeds) with surrogate $p\le0.05$
out of $216$ (9 task profiles $\times$ 24 seeds), at $K=50$ paired
surrogates per cell. The identity-swap and task-shuffle nulls preserve
the exact per-window population envelope.}
\label{tab:gateb2}
% TABLA Gateway B2: null-suite emparejada (circshift/blockshuffle/
% identityswap), K=50 por cell; FP control y readout-invariance.
% desde 04_p2/gate_b2/results/gate_b2v2_summary.json (read-only)
\begin{tabular}{llrrr}
\toprule
null type & substrate & sig MI cells & total & frac. \\
\midrule
\texttt{circshift} & LIF & 216 & 216 & 1.00 \\
\texttt{circshift} & Izhikevich & 216 & 216 & 1.00 \\
\texttt{blockshuffle} & LIF & 216 & 216 & 1.00 \\
\texttt{blockshuffle} & Izhikevich & 216 & 216 & 1.00 \\
\texttt{identityswap} & LIF & 216 & 216 & 1.00 \\
\texttt{identityswap} & Izhikevich & 216 & 216 & 1.00 \\
\texttt{taskshuf} & LIF & 216 & 216 & 1.00 \\
\texttt{taskshuf} & Izhikevich & 216 & 216 & 1.00 \\
\bottomrule
\end{tabular}

% FP control (drive_gain=0, no task signal): cells per profile =
% {"noise": {"izh": {"circshift": 0, "taskshuf": 0}, "lif": {"circshift": 0, "taskshuf": 0}}, "ramp": {"izh": {"circshift": 0, "taskshuf": 0}, "lif": {"circshift": 0, "taskshuf": 0}}, "sine": {"izh": {"circshift": 0, "taskshuf": 0}, "lif": {"circshift": 0, "taskshuf": 0}}} | criterion fp_control_clean=True.
% Readout invariance (fixed weights): 48/48 cells, p=0.0476.
\end{table}
The observed MI is above its own sample-matched surrogates in
$216/216$ pooled cells for \emph{each} of the four null types and
\emph{both} living substrates (Table~\ref{tab:gateb2}), so the effect is
not an artifact of the i.i.d.\ Poisson assumption. The task-shuffle
null is the strictest: it preserves the real count matrix exactly
(every per-neuron marginal, every cross-neuron correlation and the full
population temporal structure) and destroys only the pairing of task
signal to windows; under this strongest surrogate, both substrates still
decode the task at the maximum attainable significance. The identity-swap
surrogate is also sharp: it keeps every window's population
counts exactly and only destroys the neuron-to-task alignment, yet the
MI still collapses. The dead-drive false-positive control
(\texttt{drive\_gain}=0, no task signal, three profiles $\times$ 24
seeds $\times$ two substrates) yields \emph{zero} significant cells
($0/24$), so the surrogate test is well-calibrated at this sample size.
Finally, freezing the ridge weights learned on the real spikes and
applying them unchanged to the surrogates leaves the task MI above the
surrogate distribution in $24/24$ seed--substrate cells at the
resolution of the null sample ($p=0.0476$); because nothing is refit,
this is a statement about surviving information rather than about model
adaptability. The three criteria---null suite, false-positive control,
fixed-weight invariance---are all met (verdict \textsc{pass}).

\subsection{Gate C: does the channel carry it?}\label{sec:gatecres}
\begin{table}[H]
\centering
\caption{Gate C: canonical $\RMSE$ under the aligned channel
($\mathbf{w}_r=\mathbf{g}$) vs.\ per-seed random readouts. ``$d$'' is
$\RMSE_{\mathrm{random}}-\RMSE_{\mathrm{aligned}}$; ``paired'' is the
fraction of seeds with $d>0$. Both living substrates collapse
(3/3 seeds); the matched dead null is unaffected.}
\label{tab:gatec}
% TABLA Gateway C: canal de lectura. RMSE canon bajo wr=g (alineado) vs wr aleatorio por-seed; k_ro recalibrado.
% desde verify_wr_summary.json (sine+pulse, seeds 1/29/55)
\begin{tabular}{llrrrr}
\toprule
profile & substrate & $\widehat{\mathrm{RMSE}}_{\mathrm{aligned}}$ & $\widehat{\mathrm{RMSE}}_{\mathrm{random}}$ & $\Delta$ & $p_{\mathrm{paired}}$ \\
\midrule
\texttt{sine} & LIF & 0.183 & 0.407 & +0.224 & 3/3 \\
\texttt{sine} & Izhikevich & 0.235 & 0.396 & +0.161 & 3/3 \\
\texttt{sine} & Poisson (dead null) & 0.407 & 0.407 & +0.000 & 0/3 \\
\texttt{pulse} & LIF & 0.124 & 0.397 & +0.273 & 3/3 \\
\texttt{pulse} & Izhikevich & 0.296 & 0.361 & +0.065 & 3/3 \\
\texttt{pulse} & Poisson (dead null) & 0.397 & 0.397 & +0.000 & 0/3 \\
\bottomrule
\end{tabular}
% p_paired: fraccion de seeds con random>aligned (3 seeds);
% LIF e IZH colapsan al piso Poisson bajo wr aleatorio (d>0,
% toda seed); Poisson inmune (d=0). La informacion que queda
% utilizable solo la recupera el ridge (Gate D).
\end{table}
Randomizing the readout projection collapses \emph{both} living
substrates toward the dead floor: mean $\RMSE$ jumps from $0.12$---$0.30$
(aligned) to $0.36$---$0.41$ (randomized), in all three evaluated seeds
per profile (Table~\ref{tab:gatec}). The matched dead null is
unaffected ($\Delta=0.000$, $0/3$ seeds). Gate C therefore shows that
performance under the fixed decision map \emph{depends critically on
the alignment of the readout projection} with the input projection.
It does not by itself say the substrate is empty: a poor channel could
also fail to expose genuinely present substrate information. Gate D
settles that direction---whether the underlying spike trains retain
readout-recoverable task information---by re-learning the readout on the
same spikes.

\subsection{Gate D: the readout decides who looks like it computes}\label{sec:gatedres}
\begin{table}[H]
\centering
\caption{Gate D: same spikes, decoder sweep (pooled 5 profiles
$\times$ 6 seeds). $\dfloor=$ $\RMSE_{\mathrm{meanpd}}-
\RMSE_{\mathrm{decoder}}$ is the RMSE improvement over the
mean-predictor (no-information) floor; $\MI$ in nats;
$\MI_{\mathrm{sig}}$ is the fraction of runs with significant $\MI$.}
\label{tab:gated}
% TABLA Gateway D: reader desenchufado (MISMAS spikes, decodificadores)
% desde gate_d_summary.json (pooled 5 perfiles x 6 seeds)
\begin{tabular}{llrrrrr}
\toprule
substrate & decoder & $\widehat{\mathrm{RMSE}}$ & $\Delta\mathrm{RMSE}_{\mathrm{floor}}$ & $\mathrm{MI}\,\mathrm{(nats)}$ & $\rho$ & $\mathrm{MI}_{\mathrm{sig}}$ \\
\midrule
LIF & fixed calibrated channel & 0.168 & -0.018 & 0.651 & 0.765 & 0.800 \\
LIF & ridge (cross-validated) & 0.055 & 0.094 & 0.391 & 0.733 & 0.800 \\
Izhikevich & fixed calibrated channel & 0.239 & -0.090 & 0.385 & 0.633 & 0.800 \\
Izhikevich & ridge (cross-validated) & 0.157 & -0.008 & 0.250 & 0.458 & 0.767 \\
Poisson (dead null) & fixed calibrated channel & 0.392 & -0.243 & 0.000 & -0.000 & 0.000 \\
Poisson (dead null) & ridge (cross-validated) & 0.158 & -0.009 & 0.018 & 0.038 & 0.033 \\
\bottomrule
\end{tabular}
% $\Delta\mathrm{RMSE}_{\mathrm{floor}}$ = RMSE_floor - RMSE
% (mean-predictor floor, no information baseline; passthru ceiling=0.000)
\end{table}
Table~\ref{tab:gated}: a cross-validated ridge readout of
\emph{the same spikes} splits the substrates:
\begin{itemize}[nosep]
\item \textbf{LIF}: $\RMSE$ drops to $0.055$---the fixed channel had
left information on the table (canon-to-ridge gap $+0.113$)---and the
spikes carry the task: $\dfloor=+0.094$, $\MI$ significant in
$24/30$ runs.
\item \textbf{Izhikevich}: the fixed-channel ``tracking'' does
\emph{not} persist: ridge collapses to $0.157$, $\dfloor\approx0$
($-0.008$), $\MI$ significant in $23/30$ runs but at about a quarter
of the LIF rate. The Gate~B IZH gap was a readout-channel artifact
(confirmed by Gate~C).
\item \textbf{Dead null}: ridge $0.158\approx$ floor, $\MI\approx0$
(frac.\ $0.03$)---the matched null remains at the floor under the
cross-validated ridge readout (Hoerl & Kennard, 1970; Paninski et al., 2004).
\end{itemize}

\subsection{Gate D\,bis: does a nonlinear readout change the verdict?}\label{sec:gatedbis}
\begin{table}[H]
\centering
\caption{Gate D\,bis: same spikes and splits as Gate D, with a small
MLP residual readout on top of the cross-validated ridge. $\dfloor$
is the improvement over the mean-predictor floor; ``MI-sig'' is the
fraction of runs with significant $\MI$ (30 runs).}
\label{tab:mlp}
% TABLA Gateway D-bis: readout no-lineal (MLP 2000->32->1, L2) sobre
% las MISMAS spikes de Gate D; carried_info=meanpd_rmse-rmse_decoder.
% desde gate_d/results/mlp/mlp_summary.json
\begin{tabular}{lrrrr}
\toprule
substrate & $\dfloor$ linear & $\dfloor$ MLP & MI-sig linear & MI-sig MLP \\
\midrule
LIF & 0.094 & 0.105 & 0.800 & 0.800 \\
Izhikevich & -0.008 & 0.012 & 0.767 & 0.800 \\
Poisson (dead null) & -0.009 & -0.009 & 0.033 & 0.033 \\
\bottomrule
\end{tabular}
% MLP = ridge skip + 32-unit tanh residual net (L2, grad-clip,
% lam in {1e-3,1e-2,1e-1,1}, selected on validation).
\end{table}
A reviewer objection to linear readouts is that the best linear decoder
may hide task information that a nonlinear one would recover. We test
this directly: the MLP starts from the optimal linear readout and adds a
regularized nonlinear correction. Table~\ref{tab:mlp} shows the verdict
is essentially unchanged. LIF's $\dfloor$ moves only from $+0.094$ to
$+0.105$ (relative to the dead null, $+0.103$ to $+0.114$): the
nonlinearity is neither needed nor harmful. IZH gains a small amount
(from $\approx0$ to $+0.012$ vs.\ the null $+0.021$), still far below
LIF and consistent with the interpretation that IZH's fixed-channel
tracking is mostly a readout-channel artifact. The dead null stays at
the floor ($\dfloor=-0.009$; significant $\MI$ in $1/30$ runs). We
therefore report the linear readout as the conservative comparison and
treat the MLP as a bound, not as a new positive claim.

\subsection{Gate E: is there memory?}\label{sec:gateeres}
\begin{table}[H]
\centering
\caption{Gate E: RMSE improvement over the mean-predictor floor,
$\dfloor$\ (ridge minus floor), per profile, plus the input-only
8-tap FIR reference (positive \emph{floor} values).}
\label{tab:gatee}
% TABLA Gateway E: memoria. delta-RMSE-floor (ridge - mean-predictor) por perfil.
% desde gate_e_summary.json (6 seeds). input-FIR8 = respuesta del FIR sobre u.
\begin{tabular}{lrrrrr}
\toprule
profile & LIF & Izhikevich & Poisson & input-FIR8 \\
\midrule
\texttt{lag} & -0.139 & -0.055 & -0.016 & 0.152 \\
\texttt{integ} & -0.114 & -0.104 & -0.012 & 0.144 \\
\texttt{stair} & -0.141 & -0.025 & 0.010 & -0.016 \\
\bottomrule
\end{tabular}
% floor Gate E meanpd=0.167
\end{table}
Negative result. Three history-dependent tasks (\texttt{lag}:
4-window delayed reference; \texttt{integ}: running integral over
$8$ windows; \texttt{stair}: staircase-level switch) decoded by the
ridge on present-window spikes vs.\ an $8$-tap FIR over the input only.
$\dfloor$ is negative for every substrate
(LIF $-0.13$, IZH $-0.06$, null $-0.006$), while the input-FIR solves
the linearly recoverable ones ($+0.15$). Conclusion: no substrate shows
\emph{linearly decodable task history} in the present-window population
state. Memory is not certified here; functional memory must be
demonstrated in the closed loop, where previous activity feeds back
into the input (Carmena et al., 2003; Nicolelis, 2003).

\subsection{Gate F: recurrence ablation}\label{sec:gatef}

Section~\ref{sec:gatedres} reported that dropping recurrence
($w$ ignored) leaves the canonical RMSE unchanged. We promote this to a
gate. At $n=6$ seeds, nohub and canon agree for LIF to $<10^{-4}$
($0.1678$ vs.\ $0.1678$) and for IZH to $\approx10^{-3}$ ($0.2388$
vs.\ $0.2390$); at $n=24$ seeds the pooled gap is $-3.3\times10^{-5}$
(LIF) and $+3.5\times10^{-4}$ (IZH), i.e.\ below the resolution of the
measurement. Under the wiring and drive used here, the tested
recurrence contributes negligibly to the measured canonical RMSE, so
none of the above is a recurrence artifact. This is a statement about
\emph{this} hub at \emph{this} operating point, not a general claim
that recurrence is computationally inert.

\subsection{Architecture robustness of the ranking}\label{sec:archb}
\begin{table}[H]
\centering
\caption{Same Gate D battery (five profiles, six seeds, identical
substrates, decoders and decision map) on a second hub, architecture~B
($N=4000$, $p_{\mathrm{re}}=0.02$, variance-normalized wiring), versus
architecture~A of Sections~\ref{sec:hub}--\ref{sec:gatef}. $\dfloor$ is
the improvement over the mean-predictor floor; ``info$-$null'' subtracts
the matched dead-null floor.}
\label{tab:archb}
% TABLA arquitectura B: misma bateria Gate D (5 perfiles x 6
% seeds) sobre un segundo hub N=4000, p_re=0.02.
% desde gate_d/results_archB/archB_summary.json (A = v1 N=2000)
\begin{tabular}{llrrr}
\toprule
architecture & substrate & $\dfloor$ & info$-$null & MI-sig \\
\midrule
A ($N{=}2000$, $p_{\mathrm{re}}{=}0.01$) & LIF & 0.094 & 0.103 & 0.800 \\
A ($N{=}2000$, $p_{\mathrm{re}}{=}0.01$) & Izhikevich & -0.008 & 0.001 & 0.767 \\
B ($N{=}4000$, $p_{\mathrm{re}}{=}0.02$) & LIF & 0.101 & 0.104 & 0.800 \\
B ($N{=}4000$, $p_{\mathrm{re}}{=}0.02$) & Izhikevich & 0.000 & 0.004 & 0.800 \\
\bottomrule
\end{tabular}
% criteria_B: {"izh_weak": true, "lif_real": true, "null_dead": true, "ranking_reproduces": true}.
\end{table}
The LIF-vs-IZH gap could be an artifact of the one hub we happened to
build. To test this, we rerun the \emph{entire} Gate D battery on a
different architecture---larger and denser ($N=4000$, $p_{\mathrm{re}}=
0.02$), with the same variance-normalized wiring so the per-neuron
operating point stays comparable---while holding substrates, profiles,
seeds, decoders and the decision map fixed. The ranking reproduces
(Table~\ref{tab:archb}): LIF's carried information over the dead null is
$+0.104$ (vs.\ $+0.103$ in architecture~A) and IZH's is $+0.004$
(vs.\ $+0.001$), both far below LIF, with the dead null at the floor
($-0.004$) and the same three criteria met (LIF real, IZH weak, null
dead). The pooled canonical RMSE is likewise stable (LIF $0.165$
vs.\ $0.168$; IZH $0.242$ vs.\ $0.239$). Within the range of hubs tested,
then, the conclusion ``LIF carries linearly recoverable task information;
IZH does not'' is not bound to a specific hub size or density.

\subsection{Seed-level statistics at $n=24$}\label{sec:seeds24}
\begin{table}[H]
\centering
\caption{Seed-level paired test at $n=24$ seeds. For each seed we pool
the gain-in-$\dfloor$ of the living substrate over the dead null across
the five profiles, then run an exact paired sign-permutation over the
$24$ seeds ($10^5$ permutations).}
\label{tab:seeds24}
% TABLA seed-level n=24: efecto vs null Poisson por seed (pooled 5
% perfiles), permutation emparejada exacta.
% desde gate_d/results24/gate_d24_seed_stats.json
\begin{tabular}{lrrrr}
\toprule
substrate & $\Delta$ (vs null) & SE & paired perm $p$ & seeds same dir \\
\midrule
LIF & 0.102 & 0.001 & 1e-05 & 24/24 \\
Izhikevich & 0.006 & 0.002 & 0.01035 & 18/24 \\
\bottomrule
\end{tabular}
\end{table}
To address the small-$n$ concern directly, we reran Gate D at
$n=24$ seeds (the six canonical seeds plus eighteen new ones) with an
exact seed-level paired permutation test (Table~\ref{tab:seeds24}). LIF
is decisive: the mean effect over the null is $+0.103$ with $24/24$
seeds in the same direction ($p\approx10^{-5}$, the permutation
resolution floor). IZH is qualitatively different: a small $+0.008$
mean effect, $24/24$ seeds in the same direction
($p\approx3.6\times10^{-3}$). The pooled canonical RMSE at $n=24$ is
essentially identical to the $n=6$ battery
(LIF $0.167$ vs.\ $0.168$; IZH $0.239$ vs.\ $0.239$; ridge LIF $0.058$
vs.\ $0.055$; ridge IZH $0.153$ vs.\ $0.158$), and the dead null stays
dead (significant $\MI$ in $5\%$ of runs). We therefore report both
samplings and do not inflate the effective sample size.

\subsection{Verification}\label{sec:verify}
Four independent checks confirm the results are not implementation
artifacts:
\begin{itemize}[nosep]
\item \texttt{verify\_ridge.py}---Woodbury ridge equals a direct
least-squares solve (max $|\Delta|<4\times10^{-10}$): Gate~D numbers
are the solver's, not the optimization's;
\item \texttt{verify\_null.py}---matched dead null on two out-of-family
seeds $\{103,201\}$: $\MI=0$, $\TE=0$, $\rho=0$ exactly;
\item \texttt{verify\_wr.py}---per-seed independent random readouts
($\|\mathbf{w}^{\mathrm{rand}}\|=\|\mathbf{g}\|$, RNG $4242$): the
randomization used in Gate~C, verified against the aligned baseline
and the null;
\item \texttt{verify\_paper\_numbers.py}---every number quoted in this
v2 is re-derived from the JSON artifacts and the build fails on any
mismatch (see Section~\ref{sec:repro}).
\end{itemize}

\section{Discussion}\label{sec:discussion}

\subsection{The matched dead substrate is the falsification
anchor}\label{sec:disc1}

A scoreboard is necessary but not sufficient. Every number in
Gate~B that looks like ``computation'' is preceded by the two-line
question: \emph{does an independent Poisson process with the same
marginal rates do it too?} When it does not, the material deserves a
closer look; when it does, nothing follows about the substrate. Gate~B2
sharpens this with surrogate tests built from the substrate's own spike
train: the effect survives samples that preserve its marginals, its
short-range structure, and even its per-window population envelope, but
not their task alignment. The matched null is cheap (a rate estimate
plus a surrogate generator) and it should precede scoreboard
comparisons in the embodied-AI literature
(Kagan et al., 2022; Balci et al., 2023; Lukosevicius & Jaeger, 2009; Wringe et al., 2025).

\subsection{Readout robustness, not invariance, is the substantive test}\label{sec:disc2}

Gate~C makes the channel visible by breaking it: both living
substrates collapse when $\mathbf{w}_r$ stops matching $\mathbf{g}$.
Gate~D then asks whether a readout \emph{trained} on the same spikes
can recover the task: LIF survives ($\dfloor=+0.094$), IZH does not
($\dfloor=-0.008$). Gate~D\,bis shows a nonlinear readout does not
change the picture (LIF $+0.105$, IZH $+0.012$). Readout choice is
therefore a substantive modeling decision, not a detail. Claims like
``this material computes $\phi$'' must be re-read as ``this material's
spike statistics, \emph{via} some channel, contain readout-recoverable
task-relevant structure,'' and the channel must be shown to be
redundant---or the claim collapses to a statement about the interface,
exactly the misreading flagged in the DishBrain correspondence
(Balci et al., 2023). We use \emph{robustness} rather than
\emph{invariance} because our readouts are still a family of linear (and
one small nonlinear) models; nothing here shows invariance to arbitrary
readouts.

\subsection{What we are not claiming}\label{sec:disc3}

The measured canonical $\RMSE$ is unaffected by dropping the recurrent
connections (Gate~F); the tested recurrence contributes negligibly to
the results, and Gate E is negative, so no statement here concerns
memory. None of the numbers is a closed-loop result. A closed-loop
falsification battery is left to future work (\S\ref{sec:disc5}); this
note's protocol is the \emph{open-loop instrument} that any closed-loop
claim should first pass. Only LIF of the three materials survives the
readout controls as a carrier of task information, and only for the
rate-code tasks.

\subsection{Limits}\label{sec:disc4}

(i) The linear readout is now complemented by a small nonlinear readout
(Gate D\,bis), but not by a large network or a population model; our
nonlinear bound is deliberately weak, and larger decoders would need
more data than the $\approx66$ windows per split that the hub provides.
(ii) The matched null is \emph{single-neuron} marginal for Gate B and
sample-matched for Gate B2, including a task-shuffle that preserves the
population exactly; a null matching pairwise \emph{correlations}
across neurons would be a still stronger structure test.
(iii) Gate C randomizes one projection; joint randomization of wiring
and channel is a stronger ablation. (iv) The $6$-seed battery bounds the
exact paired permutation resolution at $1/64$; we report the $n=24$
seed-level test alongside it and keep per-seed reporting honest
(Hurlbert, 1984). (v) All results are for this hub at this operating
point; the protocol transfers, the numbers do not.

\subsection{Future work}\label{sec:disc5}

Two directions follow directly: (a) a \emph{structure}-matching null
with pairwise correlations, to test substrate \emph{structure} rather
than rates; (b) replaying the \emph{real} closed-loop data through the
same gates, converting this open-loop instrument into a closed-loop
falsification battery (Carmena et al., 2003; Zhou et al., 2026). Gate~G
covers one alternative architecture; a broader sweep (connectivity
topology, degree distribution, decision-map shape) would test the
ranking further.

\section{Reproducibility and data}\label{sec:repro}
\begin{itemize}[nosep]
\item Code: \texttt{sandbox/} git repo; Gates B/B2/C/D/D\,bis/E/F/G runners
(\texttt{gate\_b/run.py}, \texttt{04\_p2/gate\_b2/run4.py}, \texttt{gate\_d/run.py},
\texttt{gate\_d/run\_mlp.py}, \texttt{gate\_d/run\_24seeds.py},
\texttt{gate\_d/run\_archB.py}, \texttt{gate\_e/run.py}) print JSON summaries
into their \texttt{results/};
tables/figures are \emph{generated} from those JSON artifacts by
\texttt{build\_figures\_tables.py} (no hand-typed numbers).
\item Gate C: \texttt{verify/verify\_wr.py} (per-seed random readouts,
RNG $4242$) writes \texttt{verify\_wr\_summary.json}.
\item Gate D at $n=24$: \texttt{gate\_d/run\_24seeds.py} writes
\texttt{gate\_d/results24/gate\_d\_summary.json} and
\texttt{gate\_d24\_seed\_stats.json}; the $6$-seed artifacts are untouched.
Gate D\,bis: \texttt{gate\_d/run\_mlp.py} writes
\texttt{gate\_d/results/mlp/mlp\_summary.json}.
Gate~G (architecture~B): \texttt{gate\_d/run\_archB.py} patches the hub
geometry (\texttt{N}, \texttt{p\_re}) and writes
\texttt{gate\_d/results\_archB/}; the architecture~A artifacts are untouched.
\item Number audit: \texttt{verify/verify\_paper\_numbers.py} re-derives the
quoted values from the JSON artifacts and exits non-zero on any
mismatch.
\item Environment: Python \SI{3.12}{}, numpy \SI{2.5}{},
matplotlib; venv at \texttt{02\_cl/.venv312/}; \texttt{requirements.txt}
pins versions.
\item Scope: this note's claims rest on Gates B/B2/C/D/D\,bis/E/F/G only
(all artifacts above). A companion closed-loop battery lives in the
separate \texttt{surrogate\_cl/} repository; it is \emph{not}
regenerated here and none of the results in this note depend on it
(Zhou et al., 2026).
\item Build note: this v2 compiled with MiKTeX (`pdflatex`); no
Overleaf-only packages.
\end{itemize}

\section*{References}

\noindent Balci, F., Ben Hamed, S., Boraud, T., Bouret, S., Brochier, T., Brun, C., & others. (2023). A response to claims of emergent intelligence and sentience in a dish. Neuron, 111(5), 604-605.
\par
\noindent Bartolozzi, C., Indiveri, G., & Donati, E. (2022). Embodied neuromorphic intelligence. Nature Communications, 13, 1024.
\par
\noindent Brooks, R. A. (1991). Intelligence without representation. Artificial Intelligence, 47, 139-159.
\par
\noindent Carmena, J. M., & others. (2003). Learning to control a brain-machine interface for reaching and grasping by primates. PLoS Biology, 1(2), e42.
\par
\noindent D'Angelo, G., Pedersen, J. E., Hassan, T., & others. (2026). A benchmarking framework for embodied neuromorphic agents. Nature Machine Intelligence, 8, 300-312.
\par
\noindent Ernst, M. D. (2004). Permutation methods: A basis for exact inference. Statistical Science, 19(4), 676-685.
\par
\noindent Georgopoulos, A. P., Schwartz, A. B., & Kettner, R. E. (1986). Neuronal population coding of movement direction. Science, 233, 1416-1419.
\par
\noindent Gerstner, W., & Kistler, W. M. (2002). Spiking neuron models: Single neurons, populations, plasticity. Cambridge University Press.
\par
\noindent Hoerl, A. E., & Kennard, R. W. (1970). Ridge regression: Biased estimation for nonorthogonal problems. Technometrics, 12, 55-67.
\par
\noindent Hurlbert, S. H. (1984). Pseudoreplication and the design of ecological field experiments. Ecological Monographs, 54(2), 187-211.
\par
\noindent Izhikevich, E. M. (2003). Simple model of spiking neurons. IEEE Transactions on Neural Networks, 14(6), 1569-1572.
\par
\noindent Izhikevich, E. M. (2004). Which model to use for cortical spiking neurons? IEEE Transactions on Neural Networks, 15(5), 1063-1070.
\par
\noindent Jaeger, H. (2001). The echo state approach to analysing and training recurrent neural networks (GMD-Report 148). GMD.
\par
\noindent Kagan, B. J., & others. (2023). Scientific communication and the semantics of sentience. Neuron, 111(5). [Letter].
\par
\noindent Kagan, B. J., Kitchen, A. C., Tran, N. T., Habibollahi, F., Khajehnejad, M., Parker, B. J., Bhat, A., Rollo, B., Razi, A., & Friston, K. J. (2022). In vitro neurons learn and exhibit sentience when embodied in a simulated game-world. Neuron, 110(22), 3952-3969.
\par
\noindent Kraskov, A., Stoegbauer, H., & Grassberger, P. (2004). Estimating mutual information. Physical Review E, 69, 066138.
\par
\noindent Lukosevicius, M., & Jaeger, H. (2009). Reservoir computing approaches to recurrent neural network training. Computer Science Review, 3(3), 127-149.
\par
\noindent Maass, W. (1997). Networks of spiking neurons: The third generation of neural network models. Neural Networks, 10(9), 1659-1671.
\par
\noindent Maass, W., Natschlaeger, T., & Markram, H. (2002). Real-time computing without stable states: A new framework for neural computation based on perturbations. Neural Computation, 14(11), 2531-2560.
\par
\noindent Nicolelis, M. A. L. (2003). Brain-machine interfaces to restore motor function and probe neural circuits. Nature Reviews Neuroscience, 4, 417-422.
\par
\noindent Paninski, L., & others. (2004). Superlinear population encoding of dynamic hand trajectory in primary motor cortex. Journal of Neuroscience, 24(39), 8551-8561.
\par
\noindent Panzeri, S., Senatore, R., Montemurro, M. A., & Petersen, R. S. (2007). Correcting for the sampling bias in spike-train information measures. Journal of Neurophysiology, 98, 1064-1072.
\par
\noindent Pfeifer, R., & Bongard, J. (2006). How the body shapes the way we think: A new view of intelligence. MIT Press.
\par
\noindent Quiroga, R. Q., & Panzeri, S. (2009). Extracting information from neuronal populations: Information theory and decoding approaches. Nature Reviews Neuroscience, 10, 173-185.
\par
\noindent Schreiber, T. (2000). Measuring information transfer. Physical Review Letters, 85(2), 461-464.
\par
\noindent Shannon, C. E. (1948). A mathematical theory of communication. Bell System Technical Journal, 27, 379-423.
\par
\noindent Smirnova, L., Caffo, B. S., Gracias, D. H., Huang, Q., Morales Pantoja, I. E., & others. (2023). Organoid intelligence (OI): The new frontier in biocomputing and intelligence-in-a-dish. Frontiers in Science, 1, 1017235.
\par
\noindent Steinmetz, N. A., Zatka-Haas, P., Carandini, M., & Harris, K. D. (2019). Distributed coding of choice, action and engagement across the mouse brain. Nature, 576, 266-273.
\par
\noindent Theiler, J., Eubank, S., Longtin, A., Galdrikian, B., & Farmer, J. D. (1992). Testing for nonlinearity in time series: The method of surrogate data. Physica D, 58, 77-94.
\par
\noindent Walker, B. L., & Newhall, K. A. (2018). Inferring information flow in spike-train data sets using a trial-shuffle method. PLoS ONE, 13(11), e0206977.
\par
\noindent Wringe, C., Trefzer, M., & Stepney, S. (2025). Reservoir computing benchmarks: A tutorial review and critique. International Journal of Parallel, Emergent and Distributed Systems, 40(4), 313-351.
\par
\noindent Yan, M., Huang, C., Bienstman, P., Tiňo, P., Lin, W., & Sun, J. (2024). Emerging opportunities and challenges for the future of reservoir computing. Nature Communications, 15, 2056.
\par
\noindent Yik, J., Van den Berghe, K., den Blanken, D., & others. (2025). The NeuroBench framework for benchmarking neuromorphic computing algorithms and systems. Nature Communications, 16, 1545.
\par
\noindent Zhou, J., Tanneberg, D., Habibollahi, F., Loeffler, A., Lawson, K., Baccetti, V., & others. (2026). Embodied neurocomputation: A framework for interfacing biological neural cultures with scaled task-driven validation. [Preprint]. arXiv:2605.13315.
\par
\end{document}
